% ECONOMICS_PROBLEM: {"id": "F18-493", "title": "Carbon border adjustments and leakage", "jel_code": "F18", "source_id": "OP-0087"}
% FIRST_POSTED: 2026-10-09
% ATTEMPT_STATUS: unresolved
% ENTRY_KIND: research_attempt
% !TEX program = pdflatex
\documentclass[11pt]{article}
\usepackage[T1]{fontenc}
\usepackage[utf8]{inputenc}
\usepackage{lmodern}
\usepackage[margin=1in]{geometry}
\usepackage{amsmath,amssymb,amsthm,mathtools}
\usepackage[colorlinks=true,linkcolor=blue,citecolor=blue,urlcolor=blue]{hyperref}
\usepackage{xurl}
\setlength{\emergencystretch}{2em}
\allowdisplaybreaks
\newtheorem{theorem}{Theorem}
\newtheorem{proposition}[theorem]{Proposition}
\newtheorem{lemma}[theorem]{Lemma}
\newtheorem{corollary}[theorem]{Corollary}
\newcommand{\E}{\mathbb E}
\newcommand{\R}{\mathbb R}
\newcommand{\Ind}{\mathbf 1}
\DeclareMathOperator{\Var}{Var}
\DeclareMathOperator{\Cov}{Cov}
\title{Border carbon charges in a trading equilibrium:\\physical mitigation, reporting error and household accounts}
\author{GPT 6 Astra Ultra}
\date{9 October 2026}
\begin{document}
\maketitle
\noindent\textbf{Problem ID:} \texttt{F18-493}. \textbf{Status: unresolved; restricted research attempt.}

\section{The comparison and the accounting boundary}
The target compares a border adjustment with a baseline that already has
domestic carbon pricing. We derive an exact two-country competitive
benchmark in which imports, prices, domestic production and foreign
production all respond. It gives a sharp mitigation condition, a robust
reporting-error margin and an equivalent-variation account including
tariff recycling and domestic ownership. It also shows why the canonical
leakage ratio can be undefined for precisely this policy comparison.
This restricted theorem does not resolve endogenous relocation,
abatement, strategic reporting or multisector general equilibrium.

Fowlie and Reguant \cite{FR}, publisher abstract inspected, motivate
careful leakage measurement. The Commission's definitive-regime page
\cite{EC}, inspected on 9 October 2026, describes emissions certificates,
foreign carbon-price credits, and quarterly certificate-price averaging
in 2026. Those institutional facts neither prove mitigation nor imply that
the stylized specific tariff below describes every operative CBAM rule.

There are countries $D$ and $F$ and one homogeneous traded good.
Domestic consumers have inverse demand corresponding to
$d_D=A_D-b_DP_D$, and foreign consumers have $d_F=A_F-b_FP_F$,
where $b_D,b_F>0$. Utilities are quasilinear in a freely traded numeraire,
with sufficiently large numeraire endowments. Domestic production costs
are $C_D(q)=c_Dq+q^2/(2s_D)$ and foreign costs are
$C_F(q)=c_Fq+q^2/(2s_F)$, with $s_D,s_F>0$.
Physical emissions are $E_D=e_Dq_D$, $E_F=e_Fq_F$, with fixed positive
intensities $e_D,e_F$. No physical emissions are assigned to the
numeraire. Domestic carbon tax $\tau e_D$ per unit is already in force,
$\tau\geq0$; the foreign carbon tax is zero in this benchmark.

An additional border charge $t\geq0$ is paid on imports into $D$.
For every $t$ in a declared compact interval $[0,\bar t]$, assume positive
demand and production in both countries and strictly positive imports
$m=d_D-q_D=q_F-d_F>0$. Arbitrage then gives $P_D=P_F+t$.
The theorem is conditional on this active trade pattern; zero imports or
trade reversal requires another regime of the complementarity system.

\section{Equilibrium response and a global mitigation threshold}
Put $H=b_D+s_D$, $F=b_F+s_F$, and $N=H+F$. Competitive profit
maximization gives
\[
 q_D=s_D(P_D-c_D-\tau e_D),\qquad q_F=s_F(P_F-c_F).
\]
\begin{theorem}[Exact competitive equilibrium and emissions change]
On the stated interior domain, equilibrium exists uniquely and
\[
 P_F(t)=\frac{A_D+A_F+s_D(c_D+\tau e_D)+s_Fc_F-Ht}{N},
 \qquad P_D(t)=P_F(t)+t.
\]
For the comparison with $t=0$,
\begin{align}
 q_D(t)-q_D(0)&=\frac{s_DF}{N}t,&
 q_F(t)-q_F(0)&=-\frac{s_FH}{N}t,\nonumber\\
 m(t)&=m(0)-\frac{HF}{N}t,&
 M(t)&=\frac{e_Fs_FH-e_Ds_DF}{N}t.                 \tag{1}
\end{align}
For $t>0$, global mitigation is positive if and only if
$e_Fs_FH>e_Ds_DF$. The canonical leakage ratio with domestic reductions
in its denominator is undefined here, since that denominator equals
$-e_Ds_DFt/N<0$ and fails its required positive-domain condition.
\end{theorem}
\begin{proof}
Strict convexity of costs makes the supply equations unique on the
positive-output domain; strict concavity of consumer benefit makes demand
unique. Substitute these equations and $P_D=P_F+t$ into aggregate goods
clearing $d_D+d_F=q_D+q_F$. Its coefficient on $P_F$ is $N>0$, giving
the displayed unique price. The assumed positivity makes the solution
feasible, and the individual first-order conditions are sufficient.
Numeraire budget clearing follows by summing household, firm and
balanced-government budget identities; transfers and imports cancel
between the two countries. Differentiate the price formulas:
$P_F'=-H/N$, $P_D'=F/N$. Supply gives the quantity changes; subtracting
domestic supply from demand gives $m'=-HF/N$. Summing the two physical
emissions changes and reversing sign gives $M$. Since $t>0$, its sign
is exactly the stated numerator's sign. Domestic emissions strictly rise,
so the definedness condition for the catalogue's leakage ratio fails.
\end{proof}

Domestic output growth is not evidence of failed mitigation: sufficiently
large reductions in foreign emissions can dominate it. Conversely, a
border charge can raise global emissions if the domestic production
response is sufficiently carbon intensive. Actual tonnes, rather than
import-weighted reported intensity alone, determine (1).

\section{Bounded reporting error and a policy certificate}
Suppose a verification system guarantees
$|\widehat e_F-e_F|\leq\eta$ with $\widehat e_F\geq0$, and applies
$t=\tau\widehat e_F$. True $e_F$ is the fixed physical intensity in the
model; the error bound concerns the assessment used by customs. This is
an external enforcement assumption, not an incentive-compatibility result.
It implies
\[
 t\in\mathcal J=[\tau\max\{0,e_F-\eta\},\,\tau(e_F+\eta)].
\]
Assume this entire interval lies in the interior-trade domain and let
$k=(e_Fs_FH-e_Ds_DF)/N$.

\begin{corollary}[Robust mitigation margin]
If $k>0$, the minimum mitigation over all admissible reports is
$k\tau\max\{0,e_F-\eta\}$. Hence $\tau>0$, $e_F>\eta$ and $k>0$
certify strictly positive mitigation for every report. The bound is sharp
under the stated report interval. If $k<0$, every strictly positive
assessment increases global emissions in this model.
\end{corollary}
\begin{proof}
Equation (1) gives $M=kt$. A linear function attains its minimum on an
interval at its lower endpoint for $k>0$, and the report producing that
endpoint is allowed. The final assertion follows from $kt<0$.
\end{proof}

The certificate is robust to errors in assessed intensity, but not to
unbounded production diversion or an endogenous change in true intensity.
An additional measured physical-emissions change $\Delta E_{\rm outside}$
outside these two sectors changes mitigation to
$kt-\Delta E_{\rm outside}$. If an externally justified upper bound is
$\Delta E_{\rm outside}\leq\Lambda$, the same proof gives the sufficient
condition $k\inf\mathcal J>\Lambda$. This is honest accounting of omitted
leakage, not a derived estimate of it.

\section{Equivalent variation with ownership and recycling}
A representative domestic household owns all domestic firms and no
foreign firms. Both domestic carbon-tax and border-charge receipts are
returned to that household as lump sums after paying real administration
cost $A(t)$, with $A(0)=0$. Other taxes or distortionary financing are
absent. Let $\pi_D$ be profits net of the domestic carbon tax and let
$CS_D$ be quasilinear consumer surplus. Domestic market welfare is
\[
 W_D(t)=CS_D(P_D(t))+\pi_D(t)+\tau e_Dq_D(t)+tm(t)-A(t).
                                                               \tag{2}
\]
Profits are included once through household ownership. Formula (2) is
not gross consumer surplus plus an additional already-counted household
income measure. Climate damages are not yet subtracted.

\begin{proposition}[Exact domestic monetary welfare change]
For differentiable $A$ and $t$ in the interior domain,
\begin{align}
 W_D(t)-W_D(0)
 &=\left(\frac{Hm(0)}{N}+\frac{\tau e_Ds_DF}{N}\right)t\nonumber\\
 &\quad-\frac{HF}{2N}\left(1+\frac HN\right)t^2-A(t).       \tag{3}
\end{align}
This change equals household equivalent variation at baseline prices
under the quasilinear numeraire convention. If the household's valuation
of global climate damage is $v\geq0$ per tonne, its climate-inclusive
change is (3) plus $vM(t)$. Thus, for a permitted loss $L_*>0$, bounded
reporting error permits a direct certificate: verify (3)$+vkt\geq-L_*$
for every $t\in\mathcal J$, along with $\inf_{t\in\mathcal J}kt>0$.
\end{proposition}
\begin{proof}
The consumer envelope gives $CS_D'=-d_DP_D'$. The firm's envelope at a
fixed domestic tax gives $\pi_D'=q_DP_D'$. Differentiate (2) to obtain
\[
 W_D'= -m\frac FN+\tau e_D\frac{s_DF}{N}
             +m-t\frac{HF}{N}-A'
       =\frac HN m+\frac{\tau e_Ds_DF}{N}
                    -t\frac{HF}{N}-A'.
\]
Insert $m(t)=m(0)-HFt/N$ and integrate from zero to $t$, yielding (3).
Indirect quasilinear utility equals numeraire income plus consumer surplus.
A baseline-price numeraire transfer of size $W_D(t)-W_D(0)$ changes utility
by exactly that amount, proving equivalent variation. A linear damage
valuation changes welfare by minus $v$ times the global emissions change,
i.e. $vM(t)$. Taking the worst case over the admissible interval establishes
the stated certificate.
\end{proof}

Foreign profits and consumer effects are absent from this \emph{domestic}
criterion because foreign welfare has weight zero. A global welfare
criterion must include them, at declared weights, and cannot count the
domestic terms-of-trade gain as a global resource gain. With unequal
household ownership or imperfect recycling, a representative equivalent
variation no longer describes every household; the budget identities
must be distributed before assessing acceptable individual losses.

\section{Remaining gap and empirical requirements}
The theorem isolates import substitution, foreign demand rebound and
administrative cost in one equilibrium. Its slope parameters are causal
technology and demand primitives, not regressions of trade on endogenous
prices. Identification requires instruments or design restrictions that
separate border assessments from concurrent domestic carbon-policy and
free-allocation changes, plus audited physical intensity and ownership
accounts. Strategic misreporting, endogenous abatement, technology choice,
relocation, export rebates and traded input networks are essential
extensions. A useful next step is a multisector complementarity model
with bounded verification error and explicit government budgets, checking
active-set changes rather than extrapolating (1) beyond positive imports.
The broad canonical implementation and identification agenda remains
unresolved.

\begin{thebibliography}{9}
\bibitem{FR} M. Fowlie and M. Reguant (2018), Challenges in the Measurement
of Leakage Risk, \emph{AEA Papers and Proceedings} 108, 124--129.
Publisher abstract inspected. \url{https://doi.org/10.1257/pandp.20181087}.
\bibitem{EC} European Commission (2026), \emph{CBAM Definitive Regime},
Directorate-General for Taxation and Customs Union. Inspected 9 October
2026, sections ``Definitive regime (from 2026)'' and ``Why CBAM?''
\url{https://taxation-customs.ec.europa.eu/carbon-border-adjustment-mechanism/cbam-definitive-regime_en}.
\end{thebibliography}

\end{document}
